\subsection{The Taylor witness at the ramified prime}
\label{sec:ramified-witness}

Write $p=d$, $v(p)=1$, $\mathcal O=\mathcal O_{\mathbb C_p}$,
$\pi=\sqrt{-d}$ and $\rho=\gamma=1/(2(p-1))$.
We use the following local curve data. The translated coordinate $y(z)=z+O(z^2)$
has integral formal addition law, i.e.$y(z+w)\in\mathcal O[[y(z),y(w)]]$. There is an element
$T\in\mathcal O[[y]]$ satisfying the normalized endomorphism identity
of Appendix~A.1. Finally for $\delta_i\in\mathbb Z_p$ with
$\sum_i\delta_i^2=0$ we use the CM theta identity
$\prod_i\theta(\delta_i z)\in\mathcal O[[y]]$, with $\theta(0)=1$.
These are statements about one curve and its normalization.
We give the deduction of a witness at every required multidegree.

\paragraph{Derivatives in one variable.}
Let $m_n$ be the smallest valuation of $g^{(n)}(0)$ as
$g$ ranges over $\mathcal O[[y]]$, with $v(0)=+\infty$.
Only $1,y,\ldots,y^n$ contribute to an $n$th derivative, so the minimum
is attained by one of these powers and is finite since $y^n=z^n+O(z^{n+1})$.
In particular $m_0=0$. For $g\in\mathcal O[[y]]$, the addition law expands
$g(z+w)$ in the products $y(z)^i y(w)^j$ with integral coefficients.
Taking the mixed derivative of orders $a,b$ gives
$m_{a+b}\ge m_a+m_b$.
Put $N=p^j$. The rescaled parameter
$Y(z)=T(\varpi z)/\varpi$, $v(\varpi)=\rho$, satisfies
$Y\equiv\sum_{j\ge0} z^{p^j}\pmod{\pi}$ by the endomorphism recursion in
Appendix~A.1. Hence $v([z^N]T)=-(N-1)\rho$.
Legendre's formula gives $v(N!)=(N-1)/(p-1)$ and therefore
$m_N\le (N-1)\gamma$. Fixing $n>0$, write $N=an+r$ with $0\le r<n$.
Superadditivity gives
\[
 a m_n+\min_{0\le i<n}m_i\le (N-1)\gamma.
\]
Letting $j$ tend to infinity gives $m_n\le n\gamma$ for every $n$.
This argument yields a monomial $g=y^i$, $i\le n$, realizing the bound.

\paragraph{Recovering a multidegree.}
Let $L_j(z)$ be the frame's local linear forms, with
$L_*(z)=\sum_\tau u_\tau z_\tau$ one member such that $v(u_\tau)=0$
for all $\tau$. Write $\Theta=\prod_j\theta(L_j)^{c_j}$, and
$\mathcal A$ for the image of the integral formal power series ring
under substitution of the $y(L_j)$. Choose six $p$-adic integers
$\delta_0=1,\ldots,\delta_5$ with $\sum\delta_i^2=0$ as in Appendix~A.3.
Then $U(z)=\prod_{i=0}^5\Theta(\delta_i z)$ belongs to $\mathcal A$ and
has constant term one. Taking the formal inverse in the generator
variables gives $U^{-1}\in\mathcal A$.
For a multidegree $\ell$, set $n=|\ell|$ and choose $g$ as above.
The coefficient of $z^\ell$ in $g(L_*)$ is
\[
 g^{(n)}(0)\prod_\tau u_\tau^{\ell_\tau}/\prod_\tau\ell_\tau!.
\]
Its valuation is at most $\gamma n-\sum_\tau v(\ell_\tau!)$.
Take $H_0=g(L_*)U^{-1}$ and $H_i=1$ for $1\le i\le5$. Then
$\prod_i(H_i\Theta)(\delta_i z)=g(L_*(z))$.
Its indicated coefficient is a finite sum over
$\alpha_0+\cdots+\alpha_5=\ell$. The factors
$\delta_i^{|\alpha_i|}$ have nonnegative valuation when nonzero.
Consequently at least one product of coefficients of the six unscaled
series has valuation at most $\gamma n-\sum_\tau v(\ell_\tau!)$.
Six further copies contribute their constant terms one.

\paragraph{Comparison with the ramified prime sum.}
Fix the permissible pair and the scale exponent of Theorem~3.1; all estimates
below hold for all sufficiently large integer $x$ at those fixed parameters.
Write $w=a/x_1$, $h=\lceil 6b/k\rceil$ and
$B_\tau=\prod_{i=1}^h((\ell_\tau-h)_++i)$.
The square of $d\equiv q_0/2-1\pmod{q_0}$ is one modulo $q_0$,
and neither $1$ nor $q_0/2-1$ equals $-1$ modulo $q_0$.
Thus no $e\ge1$ has $-p^e\equiv1\pmod{q_0}$.
In particular $\tau$ and $-p^e\tau$ cannot both have weight $b$,
since all such places lie over $-1$. The formula for $S_p$ gives exactly
\[
 \gamma\sum_\tau\kappa_\tau-S_p
 =\tfrac12\sum_\tau\sum_{e\ge1}
       \min(\kappa_\tau p^{-e},\kappa_{-p^e\tau}).
\]
Every summand of the inner sum is at most $w$ and at most $bp^{-e}$.
Set $N=1+\lceil\log_p k\rceil$. Splitting the inner sum at $N$ bounds it
by $Nw+2b/(kp)$. The scale relations give
$Nw=o(b/(kp))$ since $N k p x/x_1\to0$. Also
$x^{42}=o(b/(kp))$ since $k=O(x^5)$ and $\log b>x^\eta$
for a fixed $\eta>0$.
For $\kappa_\tau\le\ell_\tau\le\kappa_\tau+x^{42}$ it follows that
\[
 \gamma |\ell|-S_p\le f\left(\gamma x^{42}+Nw/2+b/(kp)\right)
 \le f\big\lfloor h/p\big\rfloor
 \le\sum_\tau v(B_\tau).
\]
Indeed $h/p\ge6b/(kp)\to\infty$ and every interval of $h$
consecutive integers contains at least $\lfloor h/p\rfloor$ multiples of $p$.
The twelve coefficients just obtained therefore have valuation sum at most
\[
 S_p+\sum_\tau v(B_\tau)-\sum_\tau v(\ell_\tau!).
\]
This gives the required ramified witness with one copy of the loss budget
instead of two. The assumptions used are the stated local CM identities
and the existence of the frame with a unit row. Realizing these coefficient
products by integral adelic test forms is a separate step.

\subsection{The split and inert Taylor witnesses}
\label{sec:other-witnesses}

For these two cases set $\pi=p$ and let $q=p$ at a split prime and
$q=p^2$ at an inert prime. Take $\rho=1/(q-1)$ and
$\gamma=1/(p-1)-\rho$. Use the same one-curve identities as above,
with the following meaning for the endomorphism identity:
$T\in\mathcal O[[y]]$, $T=z+O(z^2)$, and
\[
 T(pz)=pT(z)+b_1T(z)^q+
           \sum_{i\ge2}b_iT(z)^{1+i(q-1)},\qquad
 b_i\in\mathcal O,\quad v(b_1)=0.
\]
For split $p$, this is the height-one CM parameter input, including at
$p=2$. The unscaled ring $\mathcal O[[y]]$ consists of integral
series in the local curve coordinate; its elements need not have integral
coefficients when expanded in the additive variable $z$.

\paragraph{Six copies in every residual degree.}
Choose $\varpi^{q-1}=-p/b_1$. The same coefficient induction as in
Appendix~A.1 makes $Y=T(\varpi z)/\varpi$ integral and gives
$Y-Y^q\equiv z\pmod p$. Indeed after division by $p\varpi$ the higher
terms of the endomorphism acquire factors
$\varpi^{i(q-1)}/p\in p\mathcal O$, and the recursion at degree $m>1$
divides by the unit $1-p^{m-1}$. Reduction yields
$[z^{q^j}]Y\equiv1\pmod p$.
At $N=q^j$, Legendre's formula again gives $v(N!)=(N-1)/(p-1)$.
It follows exactly as above that each multidegree $\nu$ has a six-copy
witness of valuation at most
\begin{equation}\label{eq:jet-six-general}
 \gamma|\nu|-V(\nu),\qquad V(\nu)=\sum_\tau v(\nu_\tau!).
\end{equation}
Here one uses the integral addition law, the six-theta product identity
and a form with all coefficients units. At a split prime $\gamma=0$
and $S_p=0$. Taking six more constant terms proves the entire split case.

\paragraph{Small inert primes.}
Assume $p$ inert and $p\le k$, with the integral normalized frame
and the unit leading coefficients in the inert coupling computation of
Appendix~A.3. In the permissible family $2$ splits since $d\equiv7\pmod8$.
Moreover $p$ is prime to $q_1$: the odd prime divisors of $q_0$
satisfy $d\equiv-1$ modulo each such divisor, so quadratic reciprocity
gives $(p/d)=1$ for them, and $q_1$ has the same prime divisors as $q_0$.
Set $P(\tau)=-p\tau$ on the group of places. For odd $e$ the partner is
$P^e(\tau)$. Call $(\tau,e)$ a short pair if $e$ is positive and odd,
$2e<\ell_0$, and both $\tau$ and $P^e\tau$ belong to $\Phi_b$.
Distinct short pairs have disjoint endpoints, where an endpoint may
mean either end. Two different exponents at one start would give three
members of $\Phi_b$ at the forbidden powers $0,e,e'$. Two pairs ending
at the same point with $e<e'$ would give, from the start of the longer
pair, three at powers $0,e'-e,e'$. When the exponents are equal the
starts coincide since $P$ is a permutation. If an end were a start,
there would be three members at powers $0,e,e+e'$, with
$e+e'<\ell_0$. The same reasoning excludes a loop (take $e'=e$).
Each use of sparseness is valid since $P^i=(-1)^i p^i$ and signs
are allowed independently.

For each start take $j_\tau=\lfloor\ell_\tau/p^e\rfloor$ copies of its
pair; put $j_\tau=0$ elsewhere. Let $\nu'$ record the $p^ej_\tau$
degrees at each start, $\nu''$ the $j_\tau$ at its end. Their sum is
at most $\ell$: the start bound follows by rounding, while
$j_\tau\le 2b/p\le b\le\ell_{P^e\tau}$. Disjointness prevents
double charging. The uniqueness condition of the inert coupling
calculation holds: an odd competing exponent $a\le e$ with the
same endpoint and a selected start would itself be a short pair.
Thus six copies give a witness of valuation at most
$\gamma(|\nu'|+|\nu''|)-\sum j_\tau-V(\nu')-V(\nu'')$.
For the remaining six use \eqref{eq:jet-six-general} at
$\mu=\ell-\nu'-\nu''$. Multiplication in disjoint variable blocks
adds the valuations. The factorial correction at each place is at
most $3(1+\log_2(2b))$: in Legendre's digit-sum formula for the
multinomial coefficient each of its three positive digit-sum terms
divided by $p-1$ is at most the number of digits of $\ell_\tau\le2b$.

To pay for this correction put $w=a/x_1$,
$n_0=\lceil\ell_0/2\rceil$. At an inert prime
$\sum_{e\ge1,\ e\ {\rm odd}}p^{-e}=p/(p^2-1)=\gamma$, so
\[
 \gamma\sum_\tau\kappa_\tau-S_p
 =\sum_\tau\sum_{e\ {\rm odd},\,e>0}
                 \min(\kappa_\tau p^{-e},\kappa_{P^e\tau}).
\]
For $e<n_0$ all terms at a given place are at most $w$, except
possibly its unique short pair, whose contribution is at most
$j_\tau+1$. The remaining geometric tail is at most
$2b p^{-n_0}\le2b/(kp)$. For the last inequality use
$\ell_0>27s$, $2^s=m^{22}\ge x^{11}$, and $k\le 4x^5$,
so $p^{n_0-1}\ge k$ eventually. The total excess per place after
subtracting $j_\tau$, including the factorial correction, is at most
\[
 \gamma x^{42}+3(1+\log_2(2b))+1+n_0 w+2b/(kp).
\]
The first four terms together are $o(b/(kp))$, uniformly for
$p\le k$: use $k^2=O(x^{10})$, $n_0=O(\log x)$, the
exponential growth of $x_1$ in $x^{\eta/6}$ and of $b$ in $x^\eta$.
The display is at most $\lfloor h/p\rfloor\le v(B_\tau)$ for all
sufficiently large $x$. This proves the twelve-copy bound
$S_p+\sum v(B_\tau)-V(\ell)$ uniformly in the small inert primes.

\paragraph{Large inert primes.}
Now $p>k$. We use the integral resolution of the frame at this prime:
each variable, up to a unit $s$, is an integral linear combination
$sz_\tau=\sum_i a_{\tau i}L_i(z)$. For a normalized tight frame
$u_i$ with integral entries and $L_i=s u_i^\dagger z$ this follows
by taking $a_{\tau i}=c_i u_{i,\tau}$ in the chosen embedding.
The series $G_\tau=\sum_i a_{\tau i}y(L_i)$ starts with $sz_\tau$.
Thus $(\prod_\tau G_\tau^{\mu_\tau})\Theta$ has coefficient $s^{|\mu|}$
at multidegree $\mu$. This supplies any residual degree with valuation
zero, using one copy and five constant terms.

Use the exponent-one coupling at every place. Uniqueness in the
coupling calculation follows from bijectivity of $P$. Choose
\[
 i_\tau=\left\lfloor (\ell_\tau-h)_+/p\right\rfloor,\qquad
 j_\tau=\min(i_\tau,\ell_{P\tau}).
\]
As $\ell_\tau\le2b$ and $p>k$, $i_\tau\le h$.
The degree used at a place $t$ is
$p j_t+j_{P^{-1}t}\le(\ell_t-h)_++\min(h,\ell_t)=\ell_t$.
Appendix~A.3 gives valuation at most $-\sum j_\tau$ for the first
six copies. The unit-coefficient construction in the preceding
paragraph fills the remainder. For $r=(\ell_\tau-h)_+$,
$\ell_\tau!\mid r! B_\tau$, and
$v(r!)=i_\tau+v(i_\tau!)$. Also $i_\tau!\mid h!\mid B_\tau$,
the latter by the integrality of $\binom{r+h}{h}$.
Therefore $v(\ell_\tau!)-i_\tau\le 2v(B_\tau)$.
Since $h\ge x^{42}$ eventually, $r\le\kappa_\tau$, and hence
\[
 i_\tau-j_\tau\le(\kappa_\tau/p-\kappa_{P\tau})_+ .
\]
Summing this last bound costs at most the $e=1$ part of $S_p$.
This proves the bound $S_p+2\sum v(B_\tau)-V(\ell)$, uniformly for
all $p>k$ inert.

\paragraph{Inputs to the local deduction.}
Taken together, these computations give the prescribed-degree
coefficient witnesses for every prime on the scales of Theorem~3.1.
The one-curve formal parameter and theta integrality identities must
hold in the specified normalization, and the required frames must exist
with integral coefficients, a unit row and, at inert primes, the
integral resolution and normalized coupling calculation.
Passage to integral sections and the four-block test-form pairing
still require their arithmetic realizations.
